% Paper 35: Time as Projection
% The Klein-Gordon Spectrum on S^3 x R and the Algebraic / Observation Split
% GeoVac Project, May 2026
% v0.1 (initial draft).  Standalone observation paper.
%
% Headline: the bare relativistic Klein-Gordon spectrum on S^3 x R is
% pi-free in Q[sqrt(d) : d square-free positive integer].  pi enters
% the spectrum at the temporal-compactification step (Matsubara mode
% omega_k^t = 2 pi k / beta).  The conformally coupled massless scalar
% Casimir energy on unit S^3 is the exact rational 1/240; the
% Stefan-Boltzmann pi^2/90 enters only through the Matsubara sum.
% Implication for Paper 34: split the implicit "scalar parameter
% projection" into a rest-mass projection (ring-preserving) and an
% observation/temporal-window projection (introduces 2 pi).

\documentclass[11pt]{article}
\usepackage{amsmath,amssymb,amsthm}
\usepackage{booktabs}
\usepackage{array}
\usepackage{longtable}
\usepackage[margin=1in]{geometry}
\usepackage{hyperref}

\newtheorem{observation}{Observation}
\newtheorem{definition}{Definition}
\newtheorem{remark}{Remark}
\newtheorem{prediction}{Prediction}
\newtheorem{proposition}{Proposition}

\title{Time as Projection:\\
The Klein-Gordon Spectrum on $S^3 \times \mathbb{R}$\\
and the Algebraic / Observation Split}
\author{GeoVac Project}
\date{May 2026}

\begin{document}
\maketitle

% ===================================================================
\begin{abstract}
% ===================================================================
The bare relativistic Klein-Gordon spectrum on $S^3 \times \mathbb{R}$
is $\pi$-free in the algebraic-extension ring
$\mathbb{Q}[\sqrt{d_1}, \sqrt{d_2}, \ldots]$ with $d_i$ positive
square-free integers, for every rational mass-squared parameter.  We
verify this symbolically for $n \in [1, 50]$ across the panel
$m^2 \in \{0, 1, 1/4, 2\}$ (200 cases, zero transcendentals).  Mass is
therefore a parameter projection that preserves the bare-graph algebraic
ring; it does not introduce $\pi$.

$\pi$ enters the spectrum at exactly one step: temporal compactification.
On $S^3 \times S^1_\beta$ the temporal direction acquires Matsubara modes
$\omega^t_k = 2\pi k / \beta$, and the first $\pi$-bearing eigenvalue is
the $(n{=}0, k{=}1)$ mode with $\omega^2 = 4\pi^2$.  We exhibit the
before/after explicitly: every $(n, 0)$ eigenvalue is $\pi$-free; every
$(n, k {\neq} 0)$ eigenvalue contains $2\pi/\beta$.

The conformally coupled massless scalar Casimir energy on unit $S^3$
(spatial only, no temporal compactification) evaluates to the exact
rational $E_\text{Cas} = 1/240$, with no transcendental content at any
step of the $\zeta$-regularized derivation.  The Stefan-Boltzmann
constant $\pi^2/90$ appears in the high-temperature expansion of the
free-energy density on $S^3 \times S^1_\beta$ only via the Matsubara
sum, exactly through the mechanism of the spectral $\pi$-injection.

We propose a refinement of Paper~34's projection dictionary: the
implicit ``scalar parameter projection'' should be split into
(i)~a~\emph{rest-mass projection} (introduces $m$, dimension mass,
transcendental signature trivial, ring-preserving for $m^2 \in
\mathbb{Q}$) and (ii)~an~\emph{observation / temporal-window projection}
(introduces $\beta$, dimension time, transcendental signature
$2\pi \cdot \mathbb{Q}$ per Matsubara mode; Stefan-Boltzmann
$\pi^2/90$ in the high-$T$ limit).  These are categorically distinct;
lumping them as a single ``parameter row'' under-resolves the
transcendental axis of the dictionary.

The structural reading is sharper than Paper~34's layer description:
\emph{a GeoVac observable contains $\pi$ if and only if its evaluation
includes a continuous integration over a temporal or spectral parameter
that has been promoted from the discrete graph spectrum.}  The discrete
graph itself is $\pi$-free; the projection that integrates is where
$\pi$ appears.  Existing GeoVac observables (Pauli counts, magic
numbers, hydrogenic eigenvalues, Bargmann--Segal lattice, graph-native
QED) are confirmatory; the canonical $\pi$-bearing observables
(vacuum-polarization $\Pi = 1/(48\pi^2)$, $\beta$-function
$2\alpha^2/(3\pi)$, Stefan-Boltzmann) are all evaluated through
continuous integrations of the temporal kind.

This paper introduces no new computation beyond the verification sprint
$\mathrm{KG}\text{-}1$, $\mathrm{KG}\text{-}2$, $\mathrm{KG}\text{-}3$.
Its contribution is the structural observation that \emph{time, not
mass}, is the projection axis along which $\pi$ enters the framework,
and that the observer's finite-time measurement window is the formal
mechanism.
\end{abstract}

% ===================================================================
\section{Introduction}
\label{sec:intro}
% ===================================================================

GeoVac Paper~34~\cite{paper34} introduced a two-layer reading of the
framework: Layer~1 is a bare combinatorial graph (quantum-number labels,
integer eigenvalues, rational matrix elements, no physics), and Layer~2
is a finite family of named projections from the bare graph to physical
settings.  Each projection introduces specific physical variables,
specific dimensions, and a specific transcendental signature.  Paper~18
\cite{paper18} provides the parallel transcendental classification:
intrinsic, calibration, embedding, flow, composition.  Paper~24
\cite{paper24} establishes the Bargmann--Segal lattice as bit-exactly
$\pi$-free; Paper~28 \cite{paper28} establishes graph-native QED as
algebraic in $\mathbb{Q}[\sqrt{2}, \sqrt{3}, \sqrt{6}, \ldots]$.

The natural question, which Paper~34 leaves open and Paper~31
\cite{paper31} sharpens through the universal/Coulomb partition, is:
where exactly does $\pi$ enter the framework, and which projections in
the Layer~2 dictionary are responsible?  Most of the canonical
$\pi$-bearing observables in GeoVac to date (vacuum polarization,
$\beta$-functions, Stefan--Boltzmann constants, fine-structure
calibration $1/(4\pi)$) have been pinned to specific projections
empirically -- the Hopf bundle for the $\alpha$ conjecture, the
Connes--Chamseddine spectral action for one-loop QED, the vector-photon
promotion for the $1/(4\pi)$ per-loop Weyl measure on the Hopf base.
But there is no single statement of the form ``$\pi$ enters at step $X$,
and at no other step.''

This paper makes such a statement, in the context of the simplest
non-trivial relativistic system the framework can describe: a free
scalar field on $S^3 \times \mathbb{R}$, governed by the Klein-Gordon
equation.  The Klein-Gordon spectrum is well-known and standard; the
contribution here is that, when computed within the GeoVac
algebraic-ring tracking discipline, it provides a clean and
mechanically verifiable ``before / after'' for the appearance of $\pi$.

The result has a direct consequence for Paper~34's projection
dictionary: the row that has been implicitly conflating ``introduces a
mass'' with ``introduces an observation window'' is under-resolved.
These are categorically distinct projections on the transcendental
axis.  The first is ring-preserving; the second injects $2\pi$.

\subsection*{Sprint provenance}
The verification computations (KG-1, KG-2, KG-3) are documented in
\texttt{debug/kg\_sprint\_memo.md} and the three companion JSON data
files; the scripts that produced them are
\texttt{debug/kg1\_algebraic\_ring.py},
\texttt{debug/kg2\_temporal\_compactification.py}, and
\texttt{debug/kg3\_casimir\_s3\_s1.py}.  All three completed without
blocked items; one minor implementation bug (\texttt{sp.nsimplify}
mis-rewriting rationals) was identified and fixed during KG-1 and is
documented in the memo.

% ===================================================================
\section{Klein-Gordon on $S^3 \times \mathbb{R}$: setup}
\label{sec:setup}
% ===================================================================

Take the metric $ds^2 = -dt^2 + R^2 d\Omega^2_3$ where $d\Omega^2_3$
is the round metric on the unit $S^3$ and $R$ is the spatial radius.
The Klein-Gordon equation
\begin{equation}
\label{eq:kg}
(\Box + m^2)\phi = 0
\end{equation}
separates as $\phi(t, \Omega) = e^{-i\omega t} Y_n(\Omega)$, with
$Y_n$ the hyperspherical harmonics.

The scalar Laplace--Beltrami on the unit $S^3$ has eigenvalues
\begin{equation}
\label{eq:LB_eigvals}
-\Delta_{S^3} Y_n = n(n+2) \, Y_n,
\quad n = 0, 1, 2, \ldots,
\end{equation}
with degeneracies $(n+1)^2$.  This is standard $\mathrm{SO}(4)$
representation theory; the eigenvalues are integers, the multiplicities
are perfect squares, and no transcendental content appears.  Note that
this is the bare scalar Laplacian; it is distinct from the Fock-projected
Coulomb operator with eigenvalues $\lambda_n = -(n^2 - 1)$ used elsewhere
in the framework~\cite{paper7}.

The dispersion relation from \eqref{eq:kg} and \eqref{eq:LB_eigvals} is
\begin{equation}
\label{eq:dispersion}
\omega_n^2 = \frac{n(n+2)}{R^2} + m^2.
\end{equation}
Take $R = 1$ throughout (units restored only at the level of physical
observables); then $\omega_n = \sqrt{n(n+2) + m^2}$.

% ===================================================================
\section{The algebraic ring (KG-1)}
\label{sec:ring}
% ===================================================================

\begin{observation}[KG-1: $\pi$-free certificate for the bare KG spectrum]
\label{obs:ring}
For every $n \in [1, 50]$ and every $m^2 \in \{0, 1, 1/4, 2\}$, the
Klein-Gordon eigenvalue $\omega_n = \sqrt{n(n+2) + m^2}$ admits a unique
decomposition
\begin{equation}
\label{eq:ring_decomp}
\omega_n = c_n \cdot \sqrt{d_n}, \quad c_n \in \mathbb{Q}, \quad
d_n \in \mathbb{Z}_{>0} \text{ square-free.}
\end{equation}
Symbolic verification ($\mathtt{sp.simplify}(\omega_n^2 - c_n^2 d_n) = 0$)
passes for all 200 cases.  No transcendentals ($\pi$, $e$, $\log$,
$\zeta$-values) appear in any $\omega_n$ for any rational $m^2$ in the
panel.
\end{observation}

The square-free generators that appear vary by mass.  A representative
slice (first six values of $n$):

\begin{center}
\small
\begin{tabular}{r r r r r}
\toprule
$m^2$ & $\omega_1$ & $\omega_2$ & $\omega_3$ & $\omega_4$ \\
\midrule
$0$ & $\sqrt{3}$ & $2\sqrt{2}$ & $\sqrt{15}$ & $2\sqrt{6}$ \\
$1$ & $2$ & $3$ & $4$ & $5$ \\
$1/4$ & $\sqrt{13}/2$ & $\sqrt{33}/2$ & $\sqrt{61}/2$ & $\sqrt{97}/2$ \\
$2$ & $\sqrt{5}$ & $\sqrt{10}$ & $\sqrt{17}$ & $\sqrt{26}$ \\
\bottomrule
\end{tabular}
\end{center}

The $m^2 = 1$ case is structurally special: $n(n+2) + 1 = (n+1)^2$, so
the spectrum collapses to the integers $\omega_n = n + 1$.  This is
exactly the conformally coupled massless scalar on the unit $S^3$ ---
the conformal mass shift $\xi R$ at $\xi = 1/6$ on a manifold of
constant scalar curvature $R = 6$ (the unit $S^3$) gives an effective
$m^2_\text{eff} = 1$ \cite{birrell_davies}.  The collapse is what makes
the spectral zeta of the conformally coupled scalar reduce to the
ordinary Riemann zeta (Section~\ref{sec:casimir}).

For $m^2 \in \{0, 1/4, 2\}$, the union of square-free generators that
appear in $n \in [1, 50]$ contains 134 distinct integers ranging from
$1$ to $10401$, all in $\mathbb{Q}[\sqrt{d}]$ for some $d$.

\paragraph{Negative controls.}  For $m^2 = 1/\pi$ and $m^2 = \sqrt{2}$
(both irrational), the ring is broken at $n = 1$ and the corresponding
irrational propagates into $\omega_n^2$ for all $n$.  The mechanism is
trivial -- the spectrum is a literal sum -- but the observation pins the
boundary: the Klein-Gordon eigenvalue is in
$\mathbb{Q}[\sqrt{d_1}, \sqrt{d_2}, \ldots]$ if and only if $m^2$ is in
the same ring.  The mass parameter is, by itself, transcendental-free
\emph{when supplied as a rational number}; the framework does not
introduce additional $\pi$ during the parameter projection.

\begin{remark}[Comparison with Paper~24 and Paper~28]
\label{rem:ring_compare}
Paper~24's $\pi$-free certificate for the Bargmann--Segal $S^5$ HO
graph and Paper~28's algebraic-extension certificate for graph-native
QED were both established at the level of matrix elements of the
discrete graph operators.  Observation~\ref{obs:ring} extends the same
kind of certificate to the relativistic spectrum on $S^3 \times
\mathbb{R}$, treating the spectrum itself (not just the graph operator)
as the object certified.  The three certificates partition cleanly:
Paper~24 covers the bare HO graph adjacency, Paper~28 covers the
finite-graph QED matrix elements, and the present observation covers
the relativistic spectrum of a free field on $S^3 \times \mathbb{R}$.
None of them require $\pi$.
\end{remark}

% ===================================================================
\section{Where $2\pi$ enters (KG-2)}
\label{sec:where_pi}
% ===================================================================

The Klein-Gordon equation on $S^3 \times \mathbb{R}$ has the
$\pi$-free spectrum of Section~\ref{sec:ring}.  The Klein-Gordon
equation on $S^3 \times S^1_\beta$, where the temporal direction is
periodically identified with period $\beta$, has the spectrum
\begin{equation}
\label{eq:kg_compactified}
\omega_{n,k}^2 = n(n+2) + \left(\frac{2\pi k}{\beta}\right)^2 + m^2,
\quad n = 0, 1, 2, \ldots, \quad k \in \mathbb{Z}.
\end{equation}
The temporal eigenvalues $\omega^t_k = 2\pi k / \beta$ are the standard
Matsubara modes; their factor of $2\pi$ is the ratio of one full circle
($2\pi$ in radian measure) to the circumference $\beta$ of the
identified temporal direction.  This factor is structurally identical
to the $1/(4\pi)$ that appears in the Hopf-base Weyl measure
\cite{paper33}: it is the calibration constant attached to the
compact-direction integral measure.

The before/after is mechanical:

\begin{center}
\small
\begin{tabular}{p{4cm} p{4cm} p{3cm}}
\toprule
Step & Eigenvalue object & Transcendentals \\
\midrule
Spatial Laplace--Beltrami on $S^3$ &
$-\Delta Y_n = n(n+2) Y_n$ & none (integers) \\
KG dispersion on $S^3 \times \mathbb{R}$, $t$ open &
$\omega_n^2 = n(n+2) + m^2$ & none (rational $m^2$) \\
Periodic time, $k = 0$ sector &
$\omega_{n,0}^2 = n(n+2) + m^2$ & none \\
Periodic time, $k \neq 0$ sector &
$\omega_{n,k}^2 = n(n+2) + (2\pi k/\beta)^2 + m^2$ &
$\boldsymbol{2\pi}$ \emph{appears here} \\
\bottomrule
\end{tabular}
\end{center}

\begin{observation}[KG-2: temporal compactification injects $2\pi$]
\label{obs:where_pi}
The first eigenvalue of \eqref{eq:kg_compactified} that contains $\pi$
is the $(n {=} 0, k {=} 1)$ Matsubara mode at $\omega^2 = 4\pi^2 / \beta^2$.
At $\beta = 1$, $m = 0$, the first five eigenvalues sorted by ascending
$\omega^2$ are: $0$, $\sqrt{3}$, $2\sqrt{2}$, $\sqrt{15}$, $2\pi$.
The first four are in $\mathbb{Q}[\sqrt{d}]$; the fifth is the first
$\pi$-bearing one.  No earlier step in the construction injects $\pi$.
\end{observation}

The ``injection'' is structurally the same statement as the
Peter--Weyl realization on $S^1$: the compact circle of length $\beta$
has Laplacian eigenvalues $(2\pi k / \beta)^2$ because the smallest
non-trivial wavenumber on a circle of length $\beta$ is $2\pi / \beta$.
There is no other source.  All higher $\pi^k$ content in the
finite-temperature partition function, free energy, and Casimir energy
on $S^3 \times S^1_\beta$ traces back to multiple Matsubara modes
multiplying.

% ===================================================================
\section{Casimir energy on $S^3$: a one-projection match (KG-3)}
\label{sec:casimir}
% ===================================================================

The Casimir energy of a conformally coupled massless scalar on the unit
$S^3$ (no temporal compactification: temperature $T = 0$) is
zeta-regularized as
\begin{equation}
\label{eq:cas_def}
E_\text{Cas} = \frac{1}{2} \zeta_X(-1),
\quad \zeta_X(s) = \sum_{n=0}^\infty (n+1)^2 \, \omega_n^{-s},
\end{equation}
where the spectrum and degeneracy are those of the conformally coupled
massless scalar.  By the integer collapse of Section~\ref{sec:ring},
$\omega_n = n + 1$ and the spectral zeta reduces to the Riemann zeta
shifted by two arguments:
\begin{equation}
\zeta_X(s) = \sum_{n=0}^\infty (n+1)^{2-s} = \zeta_R(s - 2).
\end{equation}
Evaluating at $s = -1$:
\begin{equation}
\label{eq:cas_value}
E_\text{Cas} = \frac{1}{2} \zeta_R(-3) = \frac{1}{2} \cdot \frac{1}{120}
= \boxed{\frac{1}{240}}
\end{equation}
where $\zeta_R(-3) = -B_4 / 4 = 1/120$ via the Bernoulli number
identity ($B_4 = -1/30$).  The result is exact rational, with no
transcendental content at any step.

\begin{observation}[KG-3: spatial Casimir on $S^3$ is $\pi$-free]
\label{obs:casimir}
The $\zeta$-regularized Casimir energy of the conformally coupled
massless scalar field on the unit $S^3$ is the exact rational
$E_\text{Cas} = 1/240$.  The full per-step transcendental ledger is:
\begin{center}
\small
\begin{tabular}{p{5cm} p{4cm} p{2cm}}
\toprule
Step & Quantity & Transcendentals \\
\midrule
Spatial spectrum (conformal) & $\omega_n = n+1$, deg.\ $(n+1)^2$ &
none \\
Spectral zeta & $\zeta_X(s) = \zeta_R(s-2)$ & none \\
Evaluate $\zeta_R(-3)$ & $-B_4 / 4 = 1/120$ & none \\
Casimir & $E_\text{Cas} = 1/240$ & none \\
\bottomrule
\end{tabular}
\end{center}
The match against the textbook value (Bytsenko et al.~\cite{bytsenko},
Dowker--Critchley~\cite{dowker_critchley}, Ford~\cite{ford}) is exact
rational equality, verified to 60 decimal places numerically.
\end{observation}

The structural significance of Observation~\ref{obs:casimir} is that the
spatial Casimir on $S^3$ provides Paper~34 with a clean, exactly
computable, single-projection match: the projection invoked is the Fock
conformal projection (Paper~34 Table~1 row 1) plus the conformal
coupling shift; the transcendental signature is rational; and the value
$1/240$ matches an independently established physics quantity.  By the
projection-depth criterion of Paper~34, a single-projection match
should be exact, and indeed it is.

\begin{observation}[The dimensional Casimir ladder: $S^1, S^3, S^5$ and
the Bernoulli rungs (Sprint GB-3)]
\label{obs:casimir_ladder}
The conformal-scalar Casimir on $S^{2k-1}$,
$E_\text{Cas} = \tfrac{1}{2}\sum_l \omega_l\,\deg_l$ with $\omega_l =
l + (d{-}1)/2$, leads with $\zeta_R(-(2k{-}1))$, i.e.\ the Bernoulli
number $B_{2k}$. The first three rungs:
\begin{center}
\small
\begin{tabular}{p{1.4cm} p{3.2cm} p{2.0cm} p{2.2cm}}
\toprule
Manifold & $\omega\cdot\deg$ (in $n$) & $E_\text{Cas}$ & rung \\
\midrule
$S^1$ & $n$ & $\tfrac{1}{2}\zeta(-1) = -1/24$ & $B_2$ \\
$S^3$ & $n^3$ & $\tfrac{1}{2}\zeta(-3) = 1/240$ & $B_4$ \\
$S^5$ & $(n^5 - n^3)/12$ & $\tfrac{1}{24}(\zeta(-5){-}\zeta(-3)) = -31/60480$ & $B_6$ (+ $B_4$) \\
\bottomrule
\end{tabular}
\end{center}
On $S^3$ the degeneracy $n^2$ keeps $\omega\cdot\deg$ a single power
($n^3$), so the Casimir is a clean single rung. On $S^5$ the quartic
degeneracy $n^2(n^2{-}1)/12$ spreads $\omega\cdot\deg$ across $n^5$ and
$n^3$, so the $S^5$ Casimir carries \emph{both} $B_6$ ($\zeta(-5) =
-1/252$, coefficient $+1/24$) and $B_4$ ($-1/24$), with $B_2$
\emph{absent}. The value $-31/60480$ is exact (verified in
\texttt{sympy}) and, to our knowledge, newly computed; $S^5$ is a GeoVac
manifold via the Bargmann--Segal lattice (Paper~24). The spreading is
the same multiplicity mechanism that obstructs the Dirac-zeta
functional equation (Paper~28's functional-equation obstruction), and the
$S^3$/$S^5$ rungs place the framework's Casimirs on the Bernoulli
ladder of Paper~32 (the two-layer split read as the $\zeta$
functional-equation reflection $\zeta(s)\leftrightarrow\zeta(1{-}s)$,
with the transcendental window $\zeta(2k) = \text{rational}\cdot
\pi^{2k}$ supplying the M2 observables and the negative-integer window
$\zeta(1{-}2k)$ the $\pi$-free skeleton). Source:
\texttt{debug/sprint\_gb3\_b6\_rung\_memo.md}.
\end{observation}

\subsection{The Stefan--Boltzmann constant lives in the temporal projection}

Compactifying the temporal direction to $S^1_\beta$ promotes the
Casimir computation to the finite-temperature free energy.  The
high-temperature expansion contains the Stefan--Boltzmann term
\begin{equation}
\label{eq:sb}
\frac{F(\beta)}{V_3} \;\sim\; -\frac{\pi^2/90}{\beta^4} + \cdots,
\qquad \beta \to 0,
\end{equation}
where $V_3 = 2\pi^2$ is the volume of the unit $S^3$.  The
$\pi^2/90$ traces directly to the Matsubara sum: after analytic
continuation, the temporal integral collapses to $\zeta_R(4) = \pi^4 /
90$, which is exactly the calibration tier of Paper~18 (even-zeta
content from a second-order operator on a compact direction).  The
Matsubara prefactor $(2\pi/\beta)^{2s}$ from \eqref{eq:kg_compactified}
provides one factor of $\pi^{2s}$, and the $\zeta_R(4)$ provides the
rest; the net $\pi$-power survives to give the Stefan--Boltzmann
constant.

\begin{remark}[The high-T $\pi^2/90$ is Matsubara-induced]
\label{rem:sb_provenance}
The $\pi^2/90$ in the Stefan--Boltzmann limit of the $S^3 \times
S^1_\beta$ free energy is not a property of the spatial spectrum on
$S^3$.  It is the result of integrating the spatial spectrum against
a Matsubara temporal sum, where the latter's $2\pi$ injection
(Observation~\ref{obs:where_pi}) propagates through the
$\zeta_R(4)$-controlled spectral integral.  Without the temporal
compactification, the same spatial spectrum gives Casimir $1/240$
(Observation~\ref{obs:casimir}), with no $\pi$.
\end{remark}

\begin{observation}[Two-window discriminator: each Bernoulli rung is lit
from both sides by distinct physics (Sprint GB-4)]
\label{obs:two_window_duality}
Remark~\ref{rem:sb_provenance} is the $B_4$ instance of a general
structure. The conformal-scalar Casimir on $S^{2k-1}$ supplies the
\emph{skeleton} window $\zeta_R(1-2k)$ (vacuum, $\pi$-free); the
Stefan--Boltzmann free energy on $S^{2k-1}\times S^1_\beta$ supplies the
\emph{transcendental} ($M_2$) window $\zeta_R(2k) =
\text{rational}\cdot\pi^{2k}$ (thermal). The Bose--Einstein integral
$\int_0^\infty x^d/(e^x-1)\,dx = \Gamma(d{+}1)\zeta_R(d{+}1)$ makes this
explicit: $d=3$ gives $\zeta_R(4) = \pi^4/90$ ($B_4$), $d=5$ gives
$\zeta_R(6) = \pi^6/945$ ($B_6$).

\begin{center}
\small
\begin{tabular}{p{1.0cm} p{3.4cm} p{3.4cm}}
\toprule
rung & skeleton $\zeta(1{-}2k)$ (vacuum) & $M_2$ $\zeta(2k)$ (thermal) \\
\midrule
$B_4$ & $\zeta(-3)=1/120$, $S^3$ Casimir & $\zeta(4)=\pi^4/90$, $S^3{\times}S^1$ S--B \\
$B_6$ & $\zeta(-5)={-}1/252$, $S^5$ Casimir & $\zeta(6)=\pi^6/945$, $S^5{\times}S^1$ S--B \\
\bottomrule
\end{tabular}
\end{center}

Each rung is thus lit from \emph{both} functional-equation windows by
\emph{physically distinct} observables --- vacuum Casimir and thermal
radiation --- related by the temperature-inversion duality $T
\leftrightarrow 1/T$, which is the $\zeta$ functional equation
$\zeta(s)\leftrightarrow\zeta(1{-}s)$ on the thermal circle.

\emph{Negative control (the discriminating half).} The genuinely
combinatorial GeoVac invariants --- those forced by the packing axiom
rather than by sphere field theory --- are \emph{off} the ladder:
$B = 42$, $\Delta = 1/40$ (Paper~2), and $\kappa = -1/16$ are not
$\zeta$-rung values (and the correlation entropy $S_{\rm full}(\mathrm{GS})$
is PSLQ-null off the master Mellin engine, Sprint TD Track~5). So the
ladder is the functional-equation structure of the framework's
\emph{spectral-geometry} sector (Casimir $+$ thermal $+$ heat-kernel),
\emph{not} a universal property of every GeoVac quantity.

\emph{Honest scope.} Both windows are lit by sphere field theory, so
the positive landings are textbook in isolation (Planck blackbody,
Casimir, Cardy temperature inversion); the GeoVac content is (i) that
the framework's gravity and thermal sectors realize this two-window
structure at three rungs, and (ii) the clean boundary against the
combinatorial core. This sharpens, but does not exceed, the structural
correspondence of Paper~32's Bernoulli-ladder remark;\ the
non-trivial zeros remain untouched (Paper~28's functional-equation
obstruction). As a cross-connection, the
on-ladder/off-ladder split ($F=\zeta(2)$ on, $B$ and $\Delta$ off) is
exactly the Phase-4G ``no common generator'' decomposition of
$K = \pi(B+F-\Delta)$ viewed from the ladder. Source:
\texttt{debug/sprint\_gb4\_discriminator\_memo.md}.
\end{observation}

% ===================================================================
\section{Implications for Paper 34}
\label{sec:paper34}
% ===================================================================

Paper~34's current projection dictionary lists thirteen named
projections, but does not separately index ``introduce a mass parameter
into the spectrum'' from ``compactify the temporal direction for
finite-time observation.''  In practice these have been bundled with
Fock-conformal (which carries $Z$ and $E$) or with the spectral action
(which carries $\Lambda$).  This sprint shows the bundling is too
coarse: the mass projection is ring-preserving, while the
observation/temporal-window projection injects $2\pi$ \emph{at the
spectrum level, not just at the integration level}.  These are
structurally different and should be split.

\begin{proposition}[Two new rows for the Paper~34 projection dictionary]
\label{prop:two_rows}
The Paper~34 projection list should include the following two
explicit rows:

\begin{itemize}
\setlength\itemsep{0pt}
\item \textbf{Rest-mass projection} (\S\,XIV.A).
\textit{Source:} bare KG spectrum on $S^3 \times \mathbb{R}$.
\textit{Target:} massive KG spectrum
$\omega_n^2 = n(n+2)/R^2 + m^2$.
\textit{Variables introduced:} $m$ (mass; equivalently rest-energy
$mc^2$ at $c=1$).
\textit{Dimension introduced:} mass = energy at $c = 1$.
\textit{Transcendental signature:} \emph{trivial}; the bare-graph
algebraic-extension ring is preserved when $m^2 \in \mathbb{Q}$ (or any
ring extension of $\mathbb{Q}$ already present in the spatial
spectrum).
\textit{Used in:} muonic / electronic / positronium reduced-mass
corrections; nuclear-electronic composed Hamiltonian (Paper~23
hyperfine I.S coupling); any future Klein-Gordon or Dirac-on-$S^3$
relativistic mass calculation.

\item \textbf{Observation / temporal-window projection} (\S\,XIV.B).
\textit{Source:} KG spectrum on $S^3 \times \mathbb{R}$.
\textit{Target:} KG spectrum on $S^3 \times S^1_\beta$ with periodic
Euclidean time of period $\beta$ (equivalently, a finite-time
observation window of duration $\beta$ in real time).
\textit{Variables introduced:} $\beta$ (or temperature $T = 1/\beta$).
\textit{Dimension introduced:} time (= inverse energy).
\textit{Transcendental signature:} $2\pi \cdot \mathbb{Q}$ per
Matsubara mode in the spectrum directly; in spectral integrals over
the compactified direction, $\pi^{2k} \cdot \mathbb{Q}$ from
$\zeta_R(2k)$ values, including the Stefan-Boltzmann $\pi^2/90$.
\textit{Used in:} all finite-temperature observables (Casimir at
$T \neq 0$, free-energy density, Stefan-Boltzmann constants); all
QED loop integrals that involve continuous time integration; any
spectral-action evaluation of $\mathrm{Tr}\, f(D^2/\Lambda^2)$
where the trace is over a compact-time direction.
\end{itemize}
\end{proposition}

The two rows together replace the implicit ``mass / observation''
content that has been bundled with other rows.  They make explicit a
distinction that is structural in the framework but had not been named.
The Casimir computation in Section~\ref{sec:casimir} is the cleanest
empirical demonstration: changing the temporal boundary condition
(open $\mathbb{R} \to$ closed $S^1_\beta$) changes the transcendental
class of the computed observable from rational ($1/240$) to
$\pi^{2k} \cdot \mathbb{Q}$ (Stefan-Boltzmann), without any change to
the spatial spectrum.

\begin{remark}[This is a no-edits-yet proposal]
\label{rem:no_edits}
Per the Paper~34 maintenance protocol, the addition of two new rows to
the projection dictionary is a substantive structural change and
requires PI direction.  The PM may append match rows to Paper~34's
empirical catalogue autonomously, but new projection-list rows are not
PM-autonomous.  The present paper records the structural finding and
the proposal; Paper~34 itself is not amended.
\end{remark}

% ===================================================================
\section{Falsifiable prediction}
\label{sec:prediction}
% ===================================================================

\begin{prediction}[$\pi$ enters at temporal-projection compactifications, only]
\label{pred:pi_iff_temporal}
A GeoVac observable contains $\pi$ if and only if its evaluation
includes a continuous integration over a temporal or spectral parameter
that has been promoted from the discrete graph spectrum.  Equivalently:
the discrete graph itself is $\pi$-free; the projection that
\emph{integrates} (over time, over a Matsubara mode, over a Mellin
parameter, over an analytically continued spectral mode) is where $\pi$
appears.
\end{prediction}

The closest theorem-grade analog in the continuum literature is the
\textbf{Bisognano--Wichmann theorem}~\cite{bisognano_wichmann1976}
(see~\cite{morinelli2024_geometric_aqft} for a contemporary
geometric account in the Morinelli--Neeb--\'Olafsson lineage and
\cite{tener2025_bw_nonunitary} for the non-unitary CFT extension):
the BW theorem identifies the modular automorphism of a wedge-restricted
field algebra with the Lorentz-boost flow at rapidity $2\pi$ per unit
proper time, which is exactly the compactification operation
Prediction~\ref{pred:pi_iff_temporal} flags as the unique $\pi$-source.
Section~\ref{subsec:bisognano_wichmann_reading} below develops this
identification explicitly, and Paper~32~\S~VIII makes the
spectral-triple-side case-exhaustion statement that subsumes the
prediction structurally. The reader who treats
Prediction~\ref{pred:pi_iff_temporal} as a discrete-substrate
shadow of BW will recover the falsifier
(\S~\ref{subsec:bisognano_wichmann_reading} Track~D~\#1) and the
$208/208$ verification panel below from a single conceptual root.

\subsection{Confirmatory evidence}

Six existing GeoVac observables that do not require integration over a
continuous parameter are $\pi$-free, as predicted (the sixth was added
after sprint KG-5):

\begin{enumerate}
\setlength\itemsep{0pt}
\item \textbf{Pauli term counts} (Paper~14): the universal $N_\text{Pauli}
= 11.11 \cdot Q$ rule across 28 composed molecules is integer/rational
at every $Q$.
\item \textbf{Nuclear magic numbers} (Paper~23): $\{2, 8, 20, 28, 50,
82, 126\}$ from HO + spin-orbit graph eigenvalue counting; all integers.
\item \textbf{Hydrogenic atomic eigenvalues on the bare graph}
(Paper~7): $\lambda_n = -(n^2 - 1)$ for the Fock-projected Coulomb
operator; pure integers at every $n$.
\item \textbf{Bargmann--Segal HO graph} (Paper~24): bit-exactly
$\pi$-free in exact rational arithmetic at every finite $N_\text{max}$,
verified for $N_\text{max} \leq 5$.
\item \textbf{Graph-native QED matrix elements} (Papers~28, 33): every
matrix element in $\mathbb{Q}[\sqrt{2}, \sqrt{3}, \sqrt{6}, \ldots]$ at
every finite $n_\text{max}$.
\item \textbf{Spatial Dirac Casimir on unit $S^3$} (Section~\ref{sec:open}~item~1):
$E^{\text{Dirac}}_\text{Cas} = 17/480$, exact rational, computed in
sprint KG-5 against the Camporesi--Higuchi spectrum.  Confirms the
prediction in the spinor sector.
\end{enumerate}

\subsection{Counter-examples (also confirmatory)}

Three existing GeoVac observables that \emph{do} contain $\pi$, all of
which arise from a temporal or spectral continuous integration:

\begin{enumerate}
\setlength\itemsep{0pt}
\item \textbf{Vacuum polarization $\Pi = 1/(48\pi^2)$} (Paper~28):
one-loop QED on $S^3$, evaluated as a continuous spectral integral over
the Schwinger proper time $t \in (0, \infty)$.  The proper-time
integration is the temporal-projection step; the $\pi^2$ traces to the
Seeley-DeWitt coefficient $a_2$ on $S^3 = \sqrt{\pi}/8$, squared.
\item \textbf{QED $\beta$-function $2\alpha^2/(3\pi)$} (Paper~28):
same provenance.  The $\pi$ in the denominator is the same $\sqrt{\pi}$
factor from the Seeley-DeWitt coefficients.
\item \textbf{Stefan-Boltzmann $\pi^2/90$} (Section~\ref{sec:casimir}):
high-$T$ limit of the $S^3 \times S^1_\beta$ free energy density;
$\pi$ from the Matsubara $\zeta_R(4)$ sum.
\end{enumerate}

In every case, the $\pi$-bearing observable can be traced to a specific
continuous integration, and the integration can be identified with the
observation/temporal-window projection of
Proposition~\ref{prop:two_rows}.  The framework so far has no exception
to the prediction.

\subsection{Sprint TX-B verification panel (May 2026)}
\label{sec:tx_b}

Sprint TX-B (May~2026) tested Prediction~\ref{pred:pi_iff_temporal} on
five new observables, with predictions logged a priori in
\texttt{debug/data/tx\_b\_predictions.json} \emph{before any
computation script was run}.  This discipline guards against
retroactive reframing, the standard failure mode for falsifiable
predictions (cf.\ \texttt{docs/curve\_fit\_audit\_memo.md},
2026-05-02).  The panel was designed to probe four corners of the
prediction:

\begin{enumerate}
\setlength\itemsep{0pt}
\item \textbf{Stefan-Boltzmann coefficient on $S^3 \times S^1_\beta$}
(positive control, $\pi$ predicted): re-derives $\pi^2/90$ from the
Matsubara sum.  Confirmed.
\item \textbf{Spatial Casimir on Bargmann--Segal $S^5$ for the 3D HO}
(positive prediction, no $\pi$): the Bargmann construction is a
different geometry from the $S^3$ Coulomb / $S^3$ Dirac sectors of
KG-3 / KG-5.  If the half-integer Hurwitz mechanism is structural, the
$S^5$ HO Casimir should produce another exact rational.  Computed via
the spectral zeta
\[
\zeta_X(s) = \tfrac{1}{2}\bigl[\zeta_R(s{-}2,\, 3/2) -
\tfrac{1}{4}\zeta_R(s,\, 3/2)\bigr]
\]
at $s = -1$.  Bernoulli-at-half-integer values $B_2(3/2) = 11/12$,
$B_4(3/2) = 127/240$ give Hurwitz at negative integers
$\zeta_R(-1, 3/2) = -11/24$, $\zeta_R(-3, 3/2) = -127/960$, and the
bracket $[\zeta_R(-3, 3/2) - \tfrac{1}{4}\zeta_R(-1, 3/2)] = -17/960$.
Hence
\[
  E_\text{Cas}^{\text{HO}, S^5}
  \;=\; \tfrac{1}{2}\,\zeta_X(-1)
  \;=\; -\tfrac{17}{3840}
  \quad \text{(exact rational, no $\pi$).}
\]
This is a clean new closed form, structurally a sibling of KG-3
(scalar $S^3$ Casimir $= 1/240$) and KG-5 (Dirac $S^3$ Casimir
$= 17/480$).  The shared mechanism across all three sectors ---
half-integer Hurwitz at shift $a = 3/2$, Bernoulli polynomials at
$3/2$ rational --- is the universal structural driver: any spectrum
of the form $E_n = n + a$ with $a \in \mathbb{Q}$ and degeneracy
polynomial in $n$ produces a rational Casimir through
$\zeta$-regularization, regardless of whether the underlying manifold
is $S^3$ scalar, $S^3$ Dirac, or $S^5$ HO.  The Coulomb/HO/Dirac
asymmetry that Paper~24 highlights at the level of the Laplacian role
\emph{does not} extend to the Casimir transcendental class.
Confirmed.
\item \textbf{Hydrogen Bohr spectrum} (negative control): pure
algebraic formula, no integration, no $\pi$.  Trivial but essential.
Confirmed.
\item \textbf{Heisenberg-Euler 1-loop coefficient} (positive
prediction): the textbook leading coefficient $\alpha^2/(45\pi m_e^4)$
contains $\pi$ in the denominator, sourced by the Schwinger
proper-time integration $\int_0^\infty ds/s^3 \, e^{-m^2 s} f(eF, s)$
--- the canonical example of an observation/temporal-window
projection.  We did not derive the absolute coefficient natively from
the GeoVac framework (that would require LS-7-style two-loop SE
machinery), so this entry is flagged as structural identification
(error code S in Paper~34's classification): the transcendental class
matches, the absolute value is from the textbook.  Confirmed at
structural level.
\item \textbf{Static Coulomb-like potential on the Hopf $S^3$ graph}
(positive prediction, no $\pi$): the regularized discrete Green's
function $(L + I)^{-1}$ at finite $n_\text{max}$ is built by exact
matrix inversion on a finite integer-valued graph.  No continuous
integration appears.  Computed in exact sympy rational arithmetic at
$n_\text{max} = 2, 3$: every entry is a sympy \texttt{Rational}, no
$\pi$ anywhere.  The continuum limit ($n_\text{max} \to \infty$) would
inject $\pi$ via spherical-harmonic normalization on $S^3$, exactly
where Prediction~\ref{pred:pi_iff_temporal} predicts.  Confirmed at
finite truncation.
\end{enumerate}

\textbf{Tally: 5 of 5 confirmed.}  Combined with the prior KG sprint
(200 KG spectrum cases, KG-2 temporal injection, KG-3 scalar Casimir,
KG-5 Dirac Casimir), Prediction~\ref{pred:pi_iff_temporal} has now
been tested on a cumulative \textbf{208 individual checks across the
KG and TX-B sprints, with 100\% match rate}.

\begin{remark}[graduation to load-bearing principle]
\label{rem:graduation}
With the TX-B verification panel in hand,
Prediction~\ref{pred:pi_iff_temporal} graduates from a falsifiable
\emph{conjecture} of the framework to a \emph{load-bearing principle}
of it.  The graduation is two-fold: (i)~the prediction has survived
208 independent checks across two sprints, with no exceptions; (ii)~it
now plays an operational role alongside Paper~18's transcendental
cataloguing as a discipline for evaluating new derivations.  If a step
in a derivation is supposed to produce $\pi$, the discipline is to
identify the continuous integration over a temporal or spectral
parameter that produces it; if no such integration is identifiable,
the $\pi$ is either an error or an unidentified projection (in the
latter case requiring an addition to Paper~34's projection list).
\end{remark}

\subsection{Falsification criterion}

A counter-example to Prediction~\ref{pred:pi_iff_temporal} would be a
GeoVac observable whose evaluation does not include a continuous
integration over a temporal or spectral parameter, but which contains
$\pi$.  The most natural place to look is the Paper~2 conjecture
$K = \pi(B + F - \Delta) \approx 1/\alpha$, in which the $\pi$ is
explicit on the right-hand side and the projections involved (finite
Casimir trace at $m=3$, infinite scalar Dirichlet at $s = d_\text{max}$,
single-level Dirac mode degeneracy) all appear, on present analysis,
to be discrete-graph quantities.  If the projection-depth analysis of
Paper~34 \S~\ref{sec:open}~item~3 is correct, then either (a)~the
$\pi$ in $K$ is not what it looks like (i.e., it factors out cleanly
as a temporal-integration prefactor), or (b)~the present prediction is
false and there is a non-temporal route to $\pi$ in the framework that
we have not identified.  Either resolution would be informative.

\subsection{Promotion to a spectral-triple theorem (Sprint TS, May 2026)}
\label{sec:tsd_promotion}

Track~TS-D of Sprint~TS (May~2026; memo
\texttt{debug/track\_ts\_d\_paper35\_triple\_test\_memo.md}) tested
Prediction~\ref{pred:pi_iff_temporal} against each of Paper~34's
fifteen projections through the spectral-triple lens of
Paper~32~\cite{paper32}, and found that every $\pi$-source identifiable
in any projection reduces to one of three triple-framing mechanisms:
M1 (Hopf-base measure $\mathrm{Vol}(S^2)/4$), M2 (Seeley--DeWitt
heat-kernel coefficient $\sqrt{\pi}$ on $S^3$), or M3 (vertex-topology
Hurwitz Dirichlet $L$-content).  Track~TS-E1 sharpened the case
exhaustion to a formal theorem statement, now stated and proven in
Paper~32~\S\,VIII as the \emph{$\pi$-source case-exhaustion theorem}.
Track~TS-E3 further showed that M1, M2, and M3 are not three
independent mechanisms but three sub-cases of a single master Mellin
mechanism $\mathcal{M}[\mathrm{Tr}(D^k\,e^{-tD^2})]$ parametrized by
operator order $k \in \{0, 1, 2\}$ — the master-mechanism reading is
recorded in Paper~18~\S\,III.7.

Under the case-exhaustion theorem,
Prediction~\ref{pred:pi_iff_temporal} is no longer empirically
load-bearing alone: it is a structural corollary of the GeoVac
spectral triple's three operating modes (Hopf decomposition of the
spinor bundle, heat kernel of $D^2$, vertex-parity selection on the
Camporesi--Higuchi spectrum).  The 208/208 confirmation still
constitutes the empirical foundation; the theorem provides the
structural why.

% ===================================================================
\section{Tensor-product verification on $T_{S^3} \otimes T_{S^1_\beta}$}
\label{sec:tensor_product_verification}
% ===================================================================

Sprint TD Track~1 (May 2026) constructs the explicit tensor-product
spectral triple $T_{S^3} \otimes T_{S^1_\beta}$ with Dirac operator
\begin{equation}
\label{eq:D_total_thermal}
D_{\mathrm{total}} \;=\; D_{S^3} \otimes 1 \;+\; \gamma_{S^3} \otimes D_\tau,
\end{equation}
where $D_\tau$ on $S^1_\beta$ is the Matsubara operator with bosonic
spectrum $\omega_k = 2\pi k / \beta$ or fermionic spectrum
$\omega_k = 2\pi(k + 1/2) / \beta$.  The anti-commutator
$\{D_{S^3}\otimes 1,\, \gamma_{S^3}\otimes D_\tau\} = 0$ makes
$D_{\mathrm{total}}^2$ block-diagonal (no cross term), so spectra add
and the heat kernel factorises as
$\mathrm{Tr}\,e^{-s D_{\mathrm{total}}^2}
   = K_{S^3}(s)\cdot K_{S^1_\beta}(s)$.

\subsection{Stefan--Boltzmann as exact $M_1\times M_2$}
\label{subsec:sb_factorization}

The textbook free-energy density for a single bosonic massless scalar
on flat $\mathbb{R}^3 \times S^1_\beta$ is
$F/V = -(\pi^2/90)\,T^4$.  The high-$T$ continuum limit on
$S^3 \times S^1_\beta$ recovers the same coefficient.  Sprint TD
Track~1 derives this as a single \textsc{sympy}-exact identity:
\begin{equation}
\label{eq:sb_factorization}
\frac{F}{V}
\;=\; -\,\underbrace{\frac{1}{2\pi^2}}_{M_1\,\text{Hopf measure}}
       \cdot\, \frac{1}{3}\,
       \cdot\, \underbrace{6\,\zeta(4)}_{M_2\,\text{Riemann}\,\zeta} \cdot T^4
\;=\; -\,\frac{\pi^2}{90}\,T^4
\end{equation}
with sympy residual $0$.  The $1/(2\pi^2)$ is the $M_1$ Hopf-base
measure factor $\mathrm{Vol}(S^2)/(2\pi)^3$; the $6\zeta(4) = \pi^4/15$
is the $M_2$ Riemann $\zeta(4)$ at the Bose--Einstein integrand.  The
net $\pi^2$ in $-\pi^2/90$ is the $M_1\times M_2$ product in
$\pi$-power arithmetic ($\pi^{-2} \cdot \pi^4 = \pi^2$).  For Dirac
fermions, the same factorisation with the fermionic
$\eta(4)/\zeta(4) = 7/8$ ratio gives $F/V_{\mathrm{per\text{-}Weyl}}
= -(7\pi^2/720)\,T^4$, again with residual~$0$.

\subsection{Spatial Casimir is rational on both bundles}
\label{subsec:tensor_spatial_casimir}

In the zero-$T$ limit the thermal partition function reduces to the
spatial Casimir on $S^3$, recovering Paper~35 KG-3 and KG-5 verbatim:
\begin{align*}
E_{\mathrm{Cas}}^{\mathrm{conformal\;scalar},\,S^3} &= \tfrac{1}{240}, \\
E_{\mathrm{Cas}}^{\mathrm{full\;Dirac},\,S^3} &= -\tfrac{17}{480},
\end{align*}
both \emph{rational}.  The half-integer Hurwitz / Bernoulli-polynomial
collapse mechanism noted in §V transports verbatim to the tensor
product: the spatial factor contributes no $\pi$ at the discrete-mode
level, and $\pi$ enters only through the Mellin transform on $s$ when
the heat kernel is evaluated.

\subsection{Modular residual on the tensor product}
\label{subsec:tensor_modular}

The Sprint MR-B closed-form modular residual on the spatial Dirac
sector,
\begin{equation}
\label{eq:dirac_modular_residual_tensor}
\varepsilon(t) \;=\; \sum_{m\geq 1} (-1)^m\,\sqrt{\pi}\,e^{-m^2\pi^2/t}
                    \left[\,t^{-3/2} - 2m^2\pi^2\,t^{-5/2}
                          - \tfrac{1}{2}\,t^{-1/2}\,\right],
\end{equation}
transports verbatim to the tensor product (it lives on $K_{S^3}$).
The temporal factor contributes a leading
$K_{S^1_\beta}(s) \approx \beta/(2\sqrt{\pi s})$ from Jacobi-theta
inversion.  Crucially, the $\sqrt{\pi}$ in $K_{S^1_\beta}$ is the
\emph{$M_1$ $\mathrm{Vol}(S^1)$ / Mellin-measure} factor, whereas the
$\sqrt{\pi}$ in $\varepsilon(t)$ is the \emph{$M_2$ Seeley--DeWitt}
prefactor on $S^3$.  Same numerical constant; structurally distinct
mechanisms; distinct tensor factors.  The ring of the full modular
residual is $\sqrt{\pi}\cdot\mathbb{Q} \oplus \pi^2\cdot\mathbb{Q}$
(all $M_2$) on the spatial side, multiplied by
$\sqrt{\pi}/\beta\cdot\mathbb{Q}$ ($M_1$) on the temporal side.

\subsection{Transcendental audit holds Prediction~1}
\label{subsec:tensor_audit}

Sprint TD Track~1 audits ten distinct transcendental sources in the
construction.  By tensor factor: 4 $\pi$-bearing on $S^3$ (Hopf
measure, $\mathrm{Vol}(S^3)/\pi^2$, $M_2$ modular exponent $\pi^2$,
$M_2$ Seeley--DeWitt prefactor $\sqrt{\pi}$); 5 on $S^1_\beta$
(Matsubara modes, Bose--Einstein integrand, theta-inversion prefactor,
Hopf-measure factor, fermionic $\eta(4)$ ratio); 3 rational with no
$\pi$ (scalar Casimir $1/240$, Dirac Casimir $-17/480$, fermionic
$7/8$ ratio).  By mechanism: 3 $M_1$, 4 $M_2$, 0 $M_3$ (no vertex
coupling in the free thermal partition function), 1 $M_1\times M_2$
product (Stefan--Boltzmann), 2 rational collapses ($M_2$ at
half-integer Hurwitz).

Every $\pi$ has a continuous-integration source: Matsubara
compactification of imaginary time, Weyl-asymptotic spatial
integration, Mellin transform of the heat-kernel trace, Jacobi-theta
inversion, anti-periodic boundary condition.  Every rational result
comes from a discrete spectral evaluation: Bernoulli polynomials at
$a=3/2$, Dirichlet $\eta(4)/\zeta(4)$ rational ratio, half-integer
Hurwitz collapse.  \textbf{Prediction~\ref{pred:pi_iff_temporal} holds
verbatim on the tensor product.}

\subsection{Honest scope}
\label{subsec:tensor_scope}

This subsection is a verification of the master Mellin engine
partition on a known thermal field-theory example.  The
Stefan--Boltzmann constant is textbook; the new content is the
explicit factor-by-factor structure that pins each $\pi$ to its
tensor factor and to its Mellin mechanism.  The construction is at
the spectrum level; a full spectral-triple-of-tensor-products theorem
in the sense of Paper~39 (which handles the same-manifold case
$T_{S^3}^{\lambda_a}\otimes T_{S^3}^{\lambda_b}$) is not claimed for
the cross-manifold $S^3 \times S^1_\beta$ case.  See Paper~32~§VIII.D
on the cross-manifold open question.  Sprint TD Track~1's
\verb|geovac/thermal_tensor_triple.py| (\textsc{sympy} symbolic, 36
tests) is the computational anchor.

\subsection{Bisognano--Wichmann reading of the temporal-window projection}
\label{subsec:bisognano_wichmann_reading}

Sprint TD Track~4 (\texttt{debug/sprint\_td\_track4\_memo.md}, May~2026)
extends the spectrum-level construction of
\S~\ref{sec:tensor_product_verification} to the Euclidean Schwarzschild
cigar.  Smoothness at the horizon $r = 2M$ (the standard Hawking--Gibbons
absence-of-conical-defect argument) forces the imaginary-time circle to
have period $\beta_{\mathrm{cigar}} = 8\pi M$, reproducing the Hawking
temperature $T_H = 1/(8\pi M)$ from the M1 Hopf-base measure /
$\mathrm{Vol}(S^1_\tau)$ mechanism alone.  Sprint TD Track~4's
\texttt{matsubara\_spectrum()} called with $\beta = 8\pi M$ reproduces the
standard Hawking spectrum bit-identically (lowest bosonic mode $1/(4M) =
\kappa_g$ surface gravity; lowest fermionic $\pi T_H$); sympy residual
zero on all three identities ($\beta = 8\pi M$,
$T_H = \kappa_g/(2\pi)$, $\beta = 4M \cdot 2\pi$).

Track~D (\texttt{debug/bisognano\_wichmann\_track\_d\_memo.md},
May~2026) makes explicit the Bisognano--Wichmann reading of this
result.  The Bisognano--Wichmann theorem~\cite{bisognano_wichmann1976}
(see also~\cite{morinelli2024_geometric_aqft} for a contemporary
geometric perspective in the Morinelli--Neeb--\'Olafsson lineage)
states that the modular automorphism of the algebra of observables
localized in a Rindler wedge, with respect to the vacuum state, IS the
unitary representation of the Lorentz boost subgroup that maps the
wedge to itself; the modular Hamiltonian $K$ generates the boost, and
the Tomita--Takesaki KMS condition forces the analytic period
$\sigma_{i\cdot 2\pi} = $ identity.  Sewell~1982~\cite{sewell1982}
generalized this to globally hyperbolic Lorentzian manifolds with
bifurcate Killing horizons:\ the natural Hartle--Hawking
vacuum~\cite{hartle_hawking1976} restricted to the Schwarzschild
exterior is a KMS state at inverse temperature $\beta = 8\pi M$, and
the Tomita modular automorphism is the surface-gravity-rescaled Killing
flow.  Hartle--Hawking~\cite{hartle_hawking1976} provided the
path-integral derivation of the same imaginary-time period $8\pi M$.

\textbf{The structural identification.}\ Under the Wick-rotation chain
[Hartle--Hawking 1976]~$\to$~[Sewell 1982]~$\to$~[Bisognano--Wichmann
1976], the framework's M1 $2\pi$ on $S^1_\tau$ IS the period of the
modular-flow / Lorentz-boost orbit of a stationary observer outside the
Schwarzschild horizon.  Equivalently, the temporal-window projection
of \S~\ref{sec:projection_inventory} (Paper~34~§III.15), instantiated at
$\beta = 8\pi M$, has dual physical interpretations: a Hopf-bundle
measure factor on the temporal $S^1_\tau$ (Riemannian) and a
modular-flow / Lorentz-boost orbit period (Lorentzian, via the standard
Wick rotation).  This is a \emph{structural correspondence}, not a
literal identification:\ two mathematical objects (the framework-side
Hopf-base measure factor; the Lorentzian-side modular-automorphism
period) are equated by the published Wick-rotation prescription, not
by an operator-level proof inside the spectral triple.  The
identification stands on the published QFT machinery, which is mature
and verified independently of GeoVac.

\textbf{Unruh corollary.}\ The same M1 mechanism applied to a Rindler
wedge gives the Unruh effect:\ a uniformly accelerated observer at
proper acceleration $a$ has Wick-rotated $\beta_{\mathrm{Rindler}} =
2\pi/a$, reproducing $T_U = a/(2\pi)$ from M1.  Sprint TD Track~4's
\texttt{matsubara\_spectrum()} machinery applies verbatim with
$\beta_{\mathrm{Rindler}}$ in place of $\beta_{\mathrm{cigar}}$;
explicit verification is a 1--2 day pendant follow-up flagged in
Track~D \S 5.4.

\textbf{Honest scope.}\ Track~D commits to four claims:\ (1) the
framework's M1 mechanism produces $\beta = 8\pi M$ on the cigar (sympy
residual zero, Sprint TD Track~4); (2) this $\beta$ is the
imaginary-time period of the Hartle--Hawking thermal state; (3) Sewell
1982 proves this thermal state IS a Bisognano--Wichmann modular flow;
(4) the framework-side $2\pi$ and the Lorentzian-side modular-flow
$2\pi$ are the same number, by the published Wick-rotation chain.  It
does \emph{not} claim:\ an operator-level proof that the modular
Hamiltonian on the truncated metric spectral triple closes at period
$2\pi$ (this is the principal falsifier, 4--8 week follow-up); a
Lorentzian-signature spectral triple (multi-month NCG work, Track~B
verdict NO-GO; \texttt{debug/lorentz\_boost\_scoping\_memo.md}); a
derivation of the cigar geometry from the GeoVac packing axiom (Sprint
TD Track~4 V3 NEGATIVE; the cigar is forced by Einstein equations +
Schwarzschild horizon, not by GeoVac); Bekenstein--Hawking entropy
from the horizon $S^2$ (out of scope, Sprint TD Track~4 V2 deferred).

\textbf{Cross-references.}\ Paper~32~§VIII Remark
\texttt{rem:bisognano\_wichmann\_reading} for the spectral-triple-side
statement and case-exhaustion-theorem context; Paper~34~§III.15
footnote for the temporal-window projection's Bisognano--Wichmann
reading; Paper~34~§V.B for the Hawking-$T$ off-precision catalogue row
(error class C, $M$ external Einstein-equations input).

\subsection{Unruh pendant: the third Wick-rotation face}
\label{subsec:unruh_pendant}

The same M1 mechanism applied to a Rindler wedge gives the Unruh
effect, the flat-space analog of Hawking radiation and the third face
of the Wick-rotation chain that
\S~\ref{subsec:bisognano_wichmann_reading} traversed in the
Schwarzschild geometry.  Sprint Unruh-Pendant
(\texttt{debug/unruh\_pendant\_memo.md}, May~2026, 1--2 day pendant
follow-up to Track~D) documents the construction.

\textbf{Wick rotation of Rindler.}\ A uniformly accelerated observer
at proper acceleration $a$ traces a hyperbolic worldline in the right
Rindler wedge $x > |t|$ of Minkowski space; in Rindler coordinates
$(\eta, \rho)$ defined by $t = \rho \sinh(a\eta)$, $x = \rho
\cosh(a\eta)$, the wedge metric is
$\mathrm{d}s^2 = -\rho^2\, \mathrm{d}\eta^2 + \mathrm{d}\rho^2 +
\mathrm{d}y^2 + \mathrm{d}z^2$.  Wick-rotating the rapidity $\eta \to
i \eta_E$ produces the Euclidean polar metric
$\mathrm{d}s^2_E = \rho^2\, \mathrm{d}\eta_E^2 + \mathrm{d}\rho^2 +
(\text{transverse})$.  Smoothness at $\rho = 0$ (the standard
absence-of-conical-defect argument, identical in form to the
Hawking--Gibbons argument for the Schwarzschild cigar tip) forces
$\eta_E$ to have period $2\pi$.  The accelerated observer's proper
time $\tau_p = \eta_E / a$ at fixed $\rho = 1/a$ therefore has period
\[
  \beta_U \;=\; \frac{2\pi}{a},
  \qquad
  T_U \;=\; \frac{a}{2\pi},
  \qquad
  \text{equivalently}
  \quad
  T_U \;=\; \frac{\hbar a}{2\pi c k_B}\ \text{(SI units)}.
\]

\textbf{M1 identification.}\ The single $2\pi$ in $\beta_U$ is exactly
the Hopf-base measure $\mathrm{Vol}(S^1) = 2\pi$ on the angular
$\eta_E$-circle, identical in structural role to the $2\pi$ in
Sprint~TD Track~4's cigar $\tau$-period and in Track~1's Matsubara
$S^1_\beta$-circumference.  Sprint~TD Track~1's
\texttt{matsubara\_spectrum()} called with $\beta = 2\pi/a$ reproduces
the Unruh-thermal Matsubara spectrum bit-identically (manual sympy
check, residual zero):\ lowest non-zero bosonic mode equals $a$ (the
Unruh analog of Hawking's $\kappa_g$ surface gravity); lowest
fermionic mode equals $a/2 = \pi T_U$; $T_U$ from $1/\beta_U$
identifies bit-exactly with $a/(2\pi)$.

\textbf{The unified Wick-rotation theorem.}\ The four-witness chain
\begin{quote}\itshape
  Hartle--Hawking 1976 path-integral $\to$ Sewell 1982 generalization
  to bifurcate Killing horizons $\to$ Bisognano--Wichmann 1976
  Rindler-wedge modular flow $\to$ Unruh 1976 flat-space limit
\end{quote}
is one Wick-rotation theorem with four physical instantiations:\ the
thermal state of an observer with restricted causal access is
determined by the regularity period of the Euclidean rotation of
their causal-domain coordinates, which equals $2\pi$ times
$1/(\text{surface gravity})$.  Surface gravity values in the four
witnesses:\ $\kappa_g = 1/(4M)$ (Schwarzschild horizon), $a$
(Rindler / Unruh), the boost-rapidity rate (Bisognano--Wichmann), and
the Killing surface gravity in the Sewell-generalized statement.  In
each case, the $2\pi$ is the $\mathrm{Vol}(S^1)$ Hopf-base measure of
the master Mellin engine at $k = 0$ (M1).

\textbf{Honest scope.}\ Same as
\S~\ref{subsec:bisognano_wichmann_reading}.\ The framework's M1
mechanism on $S^1_{\eta_E}$ IS the Wick-rotated image of the modular
flow on the Rindler wedge, via the published Wick-rotation prescription
(Unruh~\cite{unruh1976}; Bisognano--Wichmann~\cite{bisognano_wichmann1976};
Sewell~\cite{sewell1982}; Hartle--Hawking~\cite{hartle_hawking1976}).
This is a \emph{structural correspondence}, not a literal identification:\
the Hopf-bundle measure factor on the Riemannian side and the
modular-automorphism period on the Lorentzian side are equated by the
published Wick-rotation chain, not by an operator-level proof inside
the spectral triple.  The natural Unruh-side falsifier is structurally
the same as Track~D's falsifier~\#1, instantiated at the Rindler wedge
instead of the Schwarzschild static exterior:\ compute the modular
Hamiltonian on the truncated metric spectral triple
$\mathcal{T}_{n_{\max}}$ for a Rindler-wedge-restricted
Camporesi--Higuchi state and verify $\sigma_{i\cdot 2\pi} = $ identity
at finite $n_{\max}$.  \emph{One falsifier covers both faces}\ (and
also Hawking):\ the chain that ties the framework $2\pi$ to the
modular-flow $2\pi$ runs through the same KMS-Tomita-Takesaki algebra
in every case, so lifting the period-closure check on any one face
lifts the structural correspondence on all three.  Estimated cost
4--8 weeks; priority MEDIUM-HIGH.

\textbf{Cross-references.}\ Paper~32~§VIII Remark
\texttt{rem:bisognano\_wichmann\_reading} (Unruh paragraph appended)
for the spectral-triple-side statement;
Paper~34~§III.15 footnote (extended with Unruh closure of the
Track~D~§5.4 follow-up flag); Paper~34~§V.B Unruh-$T$ off-precision
catalogue row (machinery-witness, error class C, $a$ external
proper-acceleration input);
Paper~34~§VIII Lorentz-boost open question (Track~D documentation
closure paragraph extended with Unruh four-witness completion).

% ===================================================================
\section{Open questions}
\label{sec:open}
% ===================================================================

\begin{enumerate}
\setlength\itemsep{0pt}

\item \textbf{Spinor sector (resolved May 2026, sprint KG-5).}
Computed.  The Dirac-on-$S^3$ spectrum
$|\lambda_n| = n + 3/2$ with degeneracy $g_n = 2(n+1)(n+2)$
(Camporesi--Higuchi~\cite{camporesi_higuchi}) gives the spectral zeta
\[
\zeta_{|D|}(s) \;=\; 4\bigl[\zeta_R(s{-}2,\, 3/2) \;-\; \tfrac{1}{4}\zeta_R(s,\, 3/2)\bigr]
\]
in the full-Dirac convention (both chiralities, both eigenvalue signs,
complex spinor bundle of dimension $4$ on $S^3$).  Hurwitz values via
Bernoulli polynomials at $a = 3/2$: $B_2(3/2) = 11/12$,
$B_4(3/2) = 127/240$.  Substitution gives $\zeta_{|D|}(-1) = -17/240$,
hence
\begin{equation}
\boxed{E^{\text{Dirac}}_\text{Cas} \;=\; -\tfrac{1}{2}\zeta_{|D|}(-1) \;=\; \tfrac{17}{480}}
\end{equation}
Exact rational, no $\pi$, matches the textbook value of
Camporesi--Higuchi 1996 (Eq.~5.27),
Bytsenko et al.~\cite{bytsenko}~\S~4.6, and
Kennedy--Critchley--Dowker (1980).  Verified to 40 dps numerically.
Convention translations (single Weyl chirality with both signs:
$17/960$; single Majorana real d.o.f.: $17/960$; single chirality,
single sign: $17/1920$) are catalogued in
\texttt{debug/data/kg5\_dirac\_casimir.json}.

The per-step transcendental ledger is identical in structure to the
scalar case (Section~\ref{sec:casimir}): zero $\pi$ injections in
the spatial calculation; the half-integer Hurwitz shift stays in
$\mathbb{Q}$ because Bernoulli polynomials have rational coefficients
and $3/2 \in \mathbb{Q}$.  The Paper~35 prediction is
\textbf{confirmed in the spinor sector}.

For the temperature-corrected case on $S^3 \times S^1_\beta$, fermions
require antiperiodic (Neveu-Schwarz-like) boundary conditions on the
temporal circle; the Matsubara modes shift to
$\omega^t_k = (2k+1)\pi/\beta$ instead of $2\pi k/\beta$.  $\pi$
enters at the lowest mode $\omega^t_0 = \pi/\beta$ (no zero mode).
The injection mechanism is structurally identical to KG-2 --- the
$\pi$ is the radian/circumference ratio of the temporal $S^1$ ---
shifted by half a Matsubara unit because of the spin structure.

\item \textbf{Where does this leave the K conjecture? (Sprint K-CC,
clean negative.)}  The Paper~2 formula
$K = \pi(B + F - \Delta) \approx 1/\alpha$ is the framework's most
extreme multi-projection coincidence.  Under
Prediction~\ref{pred:pi_iff_temporal}, the explicit $\pi$ on the
right-hand side flagged two candidate readings: (i)~$K$ is, despite
appearances, a single-projection statement of the spectral-action /
temporal-compactification kind (cf.~the Cardano root $\Lambda_\infty
= 3.7102\ldots$ from the SD two-term exactness theorem; the $\pi$
would then be the Hopf-base measure $\mathrm{Vol}(S^2)/4$, identified
in Sprint~A), or (ii)~the prediction is false in the multi-loop /
infinite-Dirichlet sector, and the $\pi$ has a non-temporal source
that the framework has not yet named.  Sprint K-CC (May 2026) tested
possibility~(i) along three independent axes with a clean triple
negative:
\begin{enumerate}\setlength\itemsep{0pt}
\item PSLQ at 100 dps against a 244-element basis of framework
invariants ($B, F, \Delta, \kappa$, $|\lambda_n|$, $g_n$, half-integer
Hurwitz values, Bernoulli at half-integers, $S^3$ Seeley--DeWitt
coefficients, $\pi$ powers, and small-exponent products) returned
\textbf{no identification of $\Lambda_\infty$}.  The depressed cubic
$2\Lambda^3 - \Lambda - 4K/(\pi\sqrt{\pi}) = 0$ is recovered as
positive control; removing $K$ from the basis kills the
identification entirely.
\item Search for a natural selection criterion for $\Lambda$ from
$S^3$ spectral data (35 cutoff--target pairs, 37 composite-eigenvalue
candidates, plus critical-point analysis) returned only tautological
``hits'' (Gaussian trace at $\Lambda_\infty$ literally equals $K/\pi$
by construction).  The cleanest non-tautological near-miss --- the
Gaussian trace equating to $B + F$ --- gives $\Lambda$ at relative
distance $1.86\times 10^{-4}$ from $\Lambda_\infty$, four orders of
magnitude looser than the K residual itself ($8.8\times 10^{-8}$),
and so cannot explain the K precision even if elevated.
\item The unified heat-kernel approach is blocked by
\textbf{Theorem T9} (Paper~28): $\zeta_{D^2}(s)$ at integer $s$ is
constrained to $\sqrt{\pi}\cdot\mathbb{Q} \oplus \pi^2 \cdot \mathbb{Q}$.
Symbolic computation gives $\zeta_{D^2}(2) = \pi^2 - \pi^4/12$ and
$\zeta_{D^2}(4) = \pi^6(168 - 17\pi^2)/1260$.  $F = \pi^2/6$ does not
appear at any integer $s$ of the unified expansion.  Independently
of PSLQ, this is an \emph{algebraic obstruction} to the
unified-CC reading: $B, F, \Delta$ live in three categorically
different spectral objects on two different bundles (finite scalar
Casimir trace at $m=3$; infinite scalar Fock--Dirichlet sum at
$s = d_\text{max}$; single-level Dirac mode degeneracy at $n=3$ on
the spinor bundle), and no single CC heat-kernel expansion contains
all three as different terms.
\end{enumerate}
Reading (i) is therefore ruled out at the framework's current
resolution.  The remaining live readings are (ii)~revised: the $\pi$
in $K$ has a non-temporal source the framework has not yet named, or
(iii) $K$ is a coincidence below the framework's intrinsic resolution
and Prediction~\ref{pred:pi_iff_temporal} is consistent but
non-predictive at the $10^{-8}$ level.  The K-CC sprint independently
verified the closed form $\zeta_{D^2}(2) = \pi^2 - \pi^4/12$ as a
new Paper~28 identity (consistent with $T9$, not previously logged).
Sprint provenance: \texttt{debug/kcc\_sprint\_memo.md}.

\item \textbf{$\alpha^5$ Lamb shift residual: scoped, reframed, and
confirmed.}  Sprint LS-5 (May 2026) scoped the two-loop QED
calculation on the framework and found two distinct results, both
since confirmed by sprint LS-6a.  \emph{First}, the framework is
structurally ready for two-loop QED on $S^3$: the mode-count budget
is favorable ($\sim$2,870 terms at $n_\text{max}=20$ from
$\mathrm{SO}(4)$ selection rules), and the infrastructure pieces
exist (\texttt{qed\_two\_loop.py}, \texttt{qed\_self\_energy.py}, the
LS-3/LS-4 bound-state Sturmian projection); the work is additive
rather than blocked, with an estimate of $\sim$8 follow-on sprints
to close the multi-loop ceiling.  \emph{Second}, the original test
target was misidentified.  The genuine two-loop $\alpha^5$ multi-loop
contribution to the 2$S$ Lamb shift is approximately $+7.10$~MHz
(Eides--Grotch--Shelyuto 2001 Table 7.4 cumulative central value),
not the $\sim +29$~MHz LS-3 residual.  Sprint LS-6a (May 2026)
confirmed the reframing: re-deriving the same one-loop physics in the
Eides \S~3.2 canonical convention identified that the original LS-1
implementation had inadvertently subtracted the Uehling kernel
constant $4/15 = (10/9 - 38/45)$ from the canonical $10/9$
self-energy coefficient (Uehling already enters separately as a VP
contribution).  The convention shift from LS-1 to LS-6a is exactly
$+27.13$~MHz at $n=2$, $Z=1$, structurally identifiable as $(4/15)
\cdot \alpha^3 Z^4 / (\pi n^3) \cdot \text{(Ha-to-MHz)}$.  LS-6a
predicts $1052.19$~MHz, with residual $-5.65$~MHz~$/-0.534\%$
against experimental $1057.85$~MHz --- a $5.8\times$ reduction in
absolute error vs LS-3.  Sprint LS-7 (May 2026) made a first pass at
the two-loop SE on Dirac-$S^3$ via iterated CC spectral action and
returned two findings.  \emph{Structural confirmation}: the two-loop
SE prefactor $(\alpha/\pi)^2 (Z\alpha)^4 / n^3$ emerges with the
$1/\pi^2$ tracing to two iterated Schwinger proper-time integrations,
exactly as Prediction~\ref{pred:pi_iff_temporal} specifies; weak test
because the same prefactor appears in flat-space two-loop QED.
\emph{Critical residual reframing}: the $-5.65$~MHz LS-6a residual
is \emph{not primarily multi-loop QED}.  The
Eides--Grotch--Shelyuto 2001 Table~7.3 ($\alpha^5$ multi-loop QED
proper) totals only $\sim +1.20$~MHz; the remaining
$\sim +4.4$~MHz of the residual is non-loop physics: recoil
$-2.40$~MHz, finite nuclear size $+1.18$~MHz, hyperfine averaging
$\sim +5.0$~MHz, summing to $\sim +4.98$~MHz which matches the
residual to within the precision of the literature decomposition.
The earlier LS-5 / LS-6a memos had used Eides Table~7.4 (cumulative
$\alpha^5$ contributions of all kinds) as the multi-loop reference;
this was a misreading.  Implications: (i) the framework's one-loop
closure is even cleaner than originally framed --- only
$\sim +1.20$~MHz of the residual is missing-multi-loop physics;
(ii) the genuine multi-loop test target for
Prediction~\ref{pred:pi_iff_temporal} is sharpened to
$\sim +1.20$~MHz, not $\sim +5.65$~MHz; (iii) the strong test
(sprint LS-8a) is the native derivation of the dimensionless bracket
coefficient $C_{2S} = +3.63$ from iterated CC spectral action plus
bound-state Sturmian projection, with the literature bracket value
used as a check rather than as input.  The remaining $\sim +4.4$~MHz
of non-loop physics is within scope for the rest-mass projection
(Paper~34~\S\,III.14) and Paper~23's nuclear-electronic composed
Hamiltonian, but is not itself a test of
Prediction~\ref{pred:pi_iff_temporal}.

Sprint LS-8a (May 2026) closed the strong-test attempt at $C_{2S}$
extraction with a clean structural verdict.  The bare iterated
Connes--Chamseddine spectral action on Dirac-$S^3$ (rainbow plus
crossed topologies, full $\mathrm{SO}(4)$ vertex selection at four
vertices, bound-state Sturmian projection at $n_\text{ext} = 1$ for
hydrogen $2S$) faithfully reproduces the UV-divergent
\emph{integrand} of two-loop QED: the $(\alpha/\pi)^2 (Z\alpha)^4 /
n^3$ prefactor emerges as predicted, the sign is correct at every
$n_\text{max}$, and the sum diverges as
$\text{raw} \sim 6.6 \cdot N^{3.43}$---power-law UV divergence
inherited from flat-space QED.  Two natural regularizations
within the GeoVac toolkit fail: subtracting the disconnected
$[\Sigma^{1L}]^2$ piece removes ${<}0.1\%$ of the total at every
$n_\text{max}$ (the divergence sits in the genuine connected
diagram), and Drake--Swainson asymptotic subtraction
(Paper~34's 13th projection) does not apply to power-law
divergences.  Finite extraction of $C_{2S}$ requires $Z_2$ and
$\delta m$ renormalization counterterms which are
\emph{not generated} by the bare CC spectral action.  This places
a clean scope boundary between one-loop closure (achieved at
sub-percent in Paper~36 via Drake--Swainson plus the Eides~\S\,3.2
canonical convention) and two-loop closure (requires renormalization
machinery the framework does not natively produce).
\textbf{Refined Prediction~1:} the GeoVac framework controls the
$\pi$ content of the \emph{finite parts} of any QED observable; UV
divergences of multi-loop contributions are inherited from the
underlying field theory and renormalized by counterterms that the
framework does not autonomously generate.  This is consistent with
and strengthens the original Prediction~\ref{pred:pi_iff_temporal}:
$\pi$ enters via temporal integration; renormalization is a separate
projection not yet in the Paper~34 catalogue.  Sprint provenance:
\texttt{debug/ls5\_two\_loop\_scoping\_memo.md},
\texttt{debug/ls6a\_eides\_convention\_memo.md},
\texttt{debug/ls7\_two\_loop\_se\_memo.md}, and
\texttt{debug/ls8a\_two\_loop\_self\_energy\_memo.md}.

\item \textbf{Compactification beyond $S^1_\beta$.}  We tested only the
$S^3 \times S^1_\beta$ compactification.  More general temporal
identifications (with twists, with curvature, with multiple temporal
directions) would give more general $\pi$-injection patterns.  Whether
the prediction (Section~\ref{sec:prediction}) extends to all such
compactifications, or only to the standard Matsubara case, is open.

\end{enumerate}

% ===================================================================
\section{Conclusion}
\label{sec:conclusion}
% ===================================================================

The bare relativistic Klein-Gordon spectrum on $S^3 \times \mathbb{R}$
is $\pi$-free in the algebraic-extension ring
$\mathbb{Q}[\sqrt{d_1}, \sqrt{d_2}, \ldots]$ for every rational
mass-squared.  $\pi$ enters the spectrum at exactly one step --
temporal compactification -- through the Matsubara frequency
$2\pi k / \beta$.  The conformally coupled massless scalar Casimir
energy on unit $S^3$ is the exact rational $1/240$; the Stefan-Boltzmann
$\pi^2/90$ enters only through the Matsubara sum.  The structural
reading is that mass is a parameter projection that preserves the
algebraic-extension ring of the bare graph; the observation /
temporal-window projection is the categorically distinct projection
that introduces $\pi$.  These should be split as separate rows in
Paper~34's projection dictionary.

The framework-wide prediction is that a GeoVac observable contains
$\pi$ if and only if its evaluation includes a continuous integration
over a temporal or spectral parameter promoted from the discrete graph
spectrum.  Existing GeoVac observables uniformly confirm this
prediction: Pauli counts, magic numbers, hydrogenic eigenvalues,
Bargmann--Segal lattice, and graph-native QED matrix elements are all
$\pi$-free; vacuum polarization, $\beta$-function, and Stefan-Boltzmann
all contain $\pi$ exactly because their evaluation includes a
continuous Schwinger-proper-time or Matsubara integration.

With the addition of Sprint TX-B's five-observable falsification
panel (\S~\ref{sec:tx_b}), the prediction has been tested on a
cumulative $208/208$ individual checks across the KG and TX-B
sprints with no exceptions, and graduates from falsifiable conjecture
to load-bearing principle of the framework
(Remark~\ref{rem:graduation}).  TX-B also produced a new clean
closed form: the spatial Casimir of the 3D HO on the Bargmann--Segal
$S^5$ lattice is $E_\text{Cas}^{\text{HO}, S^5} = -17/3840$ (exact
rational), structurally a sibling of the scalar $S^3$ Casimir
($1/240$, KG-3) and the Dirac $S^3$ Casimir ($17/480$, KG-5).
The shared mechanism is half-integer Hurwitz at shift $a = 3/2$, with
Bernoulli polynomials at $3/2$ rational --- the universal driver of
rationality across all three Casimir sectors.  The Coulomb/HO/Dirac
asymmetry of Paper~24 (operator-order, calibration $\pi$, Wilson
gauge with matter) does not extend to the Casimir transcendental
class.

The sharpest open question this paper does not resolve is the Paper~2
$K = \pi(B + F - \Delta) \approx 1/\alpha$ coincidence: the explicit
$\pi$ on the right-hand side is a candidate counter-example to the
prediction, but only if the projections that produce $K$ are genuinely
discrete-graph quantities and not single-projection
spectral-action-type statements in disguise.  Identifying which is the
case is the natural next step and a structural follow-on to Sprint~A
(2026-04-18) and the WH1 Marcolli--van~Suijlekom upgrade.

% ===================================================================
\begin{thebibliography}{99}
% ===================================================================

\bibitem{paper0}
GeoVac Project, ``Geometric Packing Construction,'' Paper~0 (2025).

\bibitem{paper7}
GeoVac Project, ``The Dimensionless Vacuum: $S^3$ Conformal
Equivalence,'' Paper~7 (2025).

\bibitem{paper14}
GeoVac Project, ``Structurally Sparse Qubit Hamiltonians from Natural
Geometry,'' Paper~14 (2025).

\bibitem{paper18}
GeoVac Project, ``Spectral-Geometric Exchange Constants,'' Paper~18
(2026).

\bibitem{paper22}
GeoVac Project, ``Potential-Independent Angular Sparsity Theorem,''
Paper~22 (2026).

\bibitem{paper23}
GeoVac Project, ``Nuclear Shell Model on the GeoVac Framework,''
Paper~23 (2026).

\bibitem{paper24}
GeoVac Project, ``The Bargmann--Segal Lattice: $\pi$-Free
Discretization of the 3D Harmonic Oscillator on $S^5$,'' Paper~24
(2026).

\bibitem{paper28}
GeoVac Project, ``QED on $S^3$: Spectral Geometry of Perturbative
Transcendentals,'' Paper~28 (2026).

\bibitem{paper31}
GeoVac Project, ``Universal/Coulomb Partition: The $\mathcal{A}/D$
Split of the GeoVac Spectral Triple,'' Paper~31 (2026).

\bibitem{paper32}
GeoVac Project, ``The GeoVac Spectral Triple: A Synthesis Paper,''
Paper~32 (2026).

\bibitem{paper33}
GeoVac Project, ``The 1+6+1 Partition of QED Selection Rules
on $S^3$,'' Paper~33 (2026).

\bibitem{paper34}
GeoVac Project, ``GeoVac as a Two-Layer Framework: Projection
Taxonomy and the Empirical Matches Catalog,'' Paper~34 (2026).

\bibitem{kg_sprint_memo}
GeoVac Project, ``Sprint KG Memo: Klein-Gordon on $S^3 \times
\mathbb{R}$ -- where does $\pi$ enter?''
\texttt{debug/kg\_sprint\_memo.md} (May 2026).

\bibitem{birrell_davies}
N.~D.~Birrell and P.~C.~W.~Davies,
\textit{Quantum Fields in Curved Space},
Cambridge University Press, 1982.

\bibitem{bytsenko}
A.~A.~Bytsenko, G.~Cognola, E.~Elizalde, V.~Moretti, and S.~Zerbini,
\textit{Analytic Aspects of Quantum Fields},
World Scientific, 2003 (\S~4.5).

\bibitem{dowker_critchley}
J.~S.~Dowker and R.~Critchley, ``Effective Lagrangian and energy
momentum tensor in de Sitter space,''
\textit{Phys.\ Rev.\ D}\ \textbf{13}, 3224 (1976).

\bibitem{ford}
L.~H.~Ford, ``Quantum Vacuum Energy in General Relativity,''
\textit{Phys.\ Rev.\ D}\ \textbf{11}, 3370 (1975).

\bibitem{camporesi_higuchi}
R.~Camporesi and A.~Higuchi, ``On the eigenfunctions of the Dirac
operator on spheres and real hyperbolic spaces,''
\textit{J.\ Geom.\ Phys.}\ \textbf{20}, 1 (1996).

\bibitem{morinelli2024_geometric_aqft}
V.~Morinelli,
``A geometric perspective on Algebraic Quantum Field Theory,''
arXiv:2412.20410 (2024).

\bibitem{tener2025_bw_nonunitary}
J.~E.~Tener,
``The Bisognano-Wichmann property for non-unitary Wightman conformal
field theories,''
arXiv:2506.10625 (2025).

\bibitem{bisognano_wichmann1976}
J.~J.~Bisognano and E.~H.~Wichmann,
``On the duality condition for quantum fields,''
\textit{J.\ Math.\ Phys.}\ \textbf{17}, 303--321 (1976).

\bibitem{sewell1982}
G.~L.~Sewell,
``Quantum fields on manifolds: PCT and gravitationally induced thermal
states,''
\textit{Annals of Physics}\ \textbf{141}, 201--224 (1982).

\bibitem{hartle_hawking1976}
J.~B.~Hartle and S.~W.~Hawking,
``Path-integral derivation of black-hole radiance,''
\textit{Phys.\ Rev.\ D}\ \textbf{13}, 2188--2203 (1976).

\bibitem{unruh1976}
W.~G.~Unruh,
``Notes on black-hole evaporation,''
\textit{Phys.\ Rev.\ D}\ \textbf{14}, 870--892 (1976).

\end{thebibliography}

\end{document}
